\section{系统架构与协同接口}
\subsection{大组总体方案与各方式关系}
系统按“领行搬运、目标观测、运动决策、执行与验证”组织。TonyPi完成抱物路线，并作为移动跟随目标；各TurtleBot依据不同观测与环境表达完成跟随，不能将共用话题理解为共用控制器。图\ref{fig:collective-system}概括方案关系；图中四类方式是四种方案，不表示这些场景都已完成同时运行。
\begin{figure}[H]\centering
\begin{tikzpicture}[node distance=7mm and 7mm]
\node[block,text width=113mm] (leader) {TonyPi：抱物、A/B/C标记路线、左/右转与放置\\为后车提供移动目标；任务状态接口按实现分别核对};
\node[block,text width=52mm,below left=9mm and -53mm of leader] (one) {方式一：背贴AR直接跟随\\相对位姿、比例控制、门控};
\node[block,text width=52mm,right=of one] (two) {方式二：地图导航跟随\\GMapping、AMCL、轨迹回退};
\node[block,text width=52mm,below=of one] (three) {方式三：实时RGB-D\\身份关联、滚动地图与规划};
\node[block,text width=52mm,right=of three] (four) {方式四：三车纵列\\点云首车，AR ID0/1后车};
\draw[flow] (leader.south)-|(one.north);
\draw[flow] (leader.south)-|(two.north);
\draw[flow] (leader.west)--++(-5mm,0)|-(three.west);
\draw[flow] (leader.east)--++(5mm,0)|-(four.east);
\node[note,below=6mm of three.south east,anchor=north,text width=110mm,xshift=4mm] {连线表示方案关联，不表示同一控制器或同一已验收部署；\\方式四队列为TonyPi → 点云首车 → ID0后车 → ID1末车。};
\end{tikzpicture}

\caption{第5大组任务与各跟随方案的关系；方式四内部还包含前后车链路}\label{fig:collective-system}
\end{figure}
\begin{longtable}{p{19mm}p{38mm}p{37mm}p{43mm}}
\caption{四种跟随方式的共同口径对照}\\\toprule
方式 & 目标观测与坐标 & 障碍响应 & 当前证据与主要缺口\\\midrule\endfirsthead
\toprule 方式 & 目标观测与坐标 & 障碍响应 & 当前证据与主要缺口\\\midrule\endhead
一 & 前车背贴ALVAR；相机或TF车体系 & 旧版停前进/转向；新版独立门控停车 & 已有AR实验说明；启动配置与定量测试待完善\\
二 & TonyPi标记转map；AMCL与历史轨迹 & 固定地图、实时扫描、move\_base/DWA与门控 & 部分软件/基础链验证；地图质量与全程跟随待测\\
三 & RGB-D身份目标转本车odom & 滚动点云代价地图、规划与制动门控 & 已有感知/控制代码；身份模板与现场验收待补\\
四 & 首车点云观测TonyPi，后车AR ID0/1相对跟随 & 前/侧扫描净空、有界转向与搜索 & 三车联调完成直行/转弯/停车；专门障碍与重捕获未测\\\bottomrule
\end{longtable}
方式四的三车联调是一个子系统场景，不等于课程四台车采用四种独立方式的联合验收。方式一与方式四均含AR相对控制，独立性需根据完整控制链、目标选择和参赛安排说明，并由全队与教师确认。不得将“车辆数”“方式数”和“小组数”互换统计。

\subsection{全队运行阶段与启动协同}
联合试验按五个阶段组织：准备阶段固定围挡、地图、标签尺寸与车辆身份；就位阶段人工摆放物块、各车确认相机/定位及控制源；启动阶段TonyPi完成抱起，跟随车在本方式必要输入满足后启动；运行阶段逐段通过A左转和B右转，分别记录领行与后车行为；收队阶段由各车独立核对进入终点及停稳后，再判定全队完成。

各组采用的启动条件不完全相同，联调时还需要分别检查。方式一新版、方式二、方式三依赖任务许可；当前人形\code{marker_route}尚需任务桥适配。方式四新C++节点直接使用标签和扫描，尚未接入\code{/competition/mission}或同一个独立速度门控。若采用人工分时启动，应在记录中注明，不能将其称为自动任务联动。

\subsection{速度发布权、故障传播与恢复}
每台底盘只运行一个实际速度控制源；方式一旧节点输出teleop，新门控方案输出navi，而方式四默认直接发布\code{/mobile_base/commands/velocity}，绕开部分速度仲裁链。实际集成需登记唯一输出及停止机制，避免把终止一个上游节点当作停止所有底盘控制器。

领行相机掉线、掉物、目标丢失、扫描/TF过期与导航失败分别记录。共享任务桥暂停不能默认停止未订阅该桥的方式四节点；全队停止需要覆盖每台车的实际命令入口。恢复前需重新确认目标身份、里程计/地图连续性和足够观察数据，不能让后车沿旧速度追赶。不同方式的原地等待、原地搜索或有界盲行策略存在差异，正文分别说明，不统一写成“目标丢失立即停车”。

\subsection{终点收队与共同完成判据}
TonyPi放置只是领行任务的结束条件。全队完成还应同时满足参赛车辆进入终点区域、速度低于约定阈值并持续停稳、未发生未恢复的目标丢失或碰撞。可写为
\begin{equation}
G_{\rm team}=G_{\rm carry}\land G_{\rm place}\land
\bigwedge_{i\in\mathcal R}(G_{{\rm arrival},i}\land G_{{\rm stopped},i})
\land G_{\rm evidence}.
\end{equation}
$\mathcal R$为本次登记的参赛车辆集合，$G_{\rm evidence}$表示记录足以核对相应结论；该式用于说明整体完成条件，实际需要结合各车运行记录检查。终点区域和实际停稳阈值需现场明确，不能直接套用某一方式的目标距离。

\subsection{感知、决策与执行}
我们将系统分为感知、决策和执行三部分。相机或深度传感器提供目标与环境信息，控制程序选择动作或导航目标，舵机和底盘完成运动。人形机器人负责物块搬运与领行，各TurtleBot负责自身的跟随控制。调试时，我们先分别检查每一条链路，再检查它们之间的配合。
\begin{figure}[H]\centering
\begin{tikzpicture}[flow/.style={-{Stealth[length=2mm]},thick,draw=Navy!85}, stage/.style={draw=Navy!45,rounded corners=2pt,fill=Navy!3,align=center,text width=35mm,minimum height=17mm,font=\small,inner sep=3mm}]
\node[stage] (v) at (0,0) {TonyPi 相机\\固定标记／道路观测};
\node[stage] (s) at (5,0) {路线状态机\\抓取、直行、转弯、放置};
\node[stage] (a) at (10,0) {动作组执行\\抱物上肢固定};
\draw[flow] (v)--(s);\draw[flow] (s)--(a);
\node[stage] (m) at (0,-2.7) {TurtleBot 相机\\背部 AR 目标};
\node[stage] (t) at (5,-2.7) {按时间戳变换至 map\\历史轨迹与回退点};
\node[stage] (n) at (10,-2.7) {\texttt{move\_base} / DWA\\地图与实时扫描};
\draw[flow] (m)--(t);\draw[flow] (t)--(n);
\node[stage] (g) at (10,-5.1) {许可与速度门控\\输出至底盘};
\draw[flow] (n)--(g);
\node[draw=Teal!55,rounded corners=2pt,fill=Teal!4,align=left,text width=85mm,inner sep=3mm,font=\footnotesize] (h) at (2.5,-5.1) {\textbf{\color{Navy}任务状态接口}\\[1.5mm]通用 \texttt{competition}：已支持心跳桥。\\最新 \texttt{marker\_route}：需适配任务状态输出。};
\draw[flow,dashed,draw=Teal] (h.east)--(g.west);
\end{tikzpicture}

\caption{人形机器人与地图跟随的数据链}\label{fig:architecture}
\end{figure}
图中的实线表示已经定义的感知和执行流程，虚线表示任务状态接口。后车能够看见前车，并不一定说明前车已经允许继续运行；例如TonyPi暂停或出现抓取问题时，跟随程序还需要相应停止。因此视觉观测、任务状态和运动输出应分别检查。

\subsection{软硬件环境与功能包}
TonyPi使用Python 3、OpenCV、NumPy、PyYAML以及设备上的舵机SDK。不同机器人的系统镜像和接口有所不同，我们在更换设备时需要检查串口、动作路径和相机输入。TurtleBot采用课程虚拟机中的Ubuntu 16.04、ROS Kinetic和Python 2.7环境，各功能包放在\code{~/catkin_ws/src}下编译。

AR识别和速度门控放在\code{ar_marker_following}中，地图跟随放在\code{multi_robot_following}中，共用的轨迹与检查函数由\code{competition_core}提供。这样可以复用功能，但也要求编译时同时包含相应依赖。各组调整程序后，还需要确认launch中填写的包名和实际脚本位置一致。

\begin{table}[H]\centering\small
\caption{多机联调中的主要接口}
\begin{tabularx}{\linewidth}{p{37mm}p{40mm}Y}\toprule
接口 & 内容 & 调试时检查的事项\\\midrule
\code{/ar_pose_marker} & 带ID和时间戳的标记位姿 & 标签类型、实际尺寸、相机内参与坐标系\\
\code{target_pose} & 转换后的目标位姿 & 目标来源与采样时刻，不能把旧位置当作新位置\\
\code{/scan} & 障碍扫描 & 有效距离、视野范围、缺测及刷新情况\\
\code{map -> odom} & 地图定位变换 & 建图与AMCL定位分别使用，避免重复发布\\
\code{/competition/mission} & 运行阶段 & 抱物、放置、完成和暂停状态的传递\\
\code{/cmd_vel_mux/input/navi} & 导航速度 & 确认实际控制程序，避免与键盘遥控争用\\
\code{/mobile_base/commands/velocity} & 方式四底盘速度 & 直接速度接口的停止与控制源检查\\\bottomrule
\end{tabularx}\end{table}

\subsection{队列关系与通信配置}
队列中的直接前车和采用的跟随方法是两个问题。方式二可以观察TonyPi的AR标记，也可以接收前车地图位姿；方式四则由首车跟踪TonyPi，后续车辆识别直接前车的标签。每轮试验都应先确定这一关系，再设置目标ID和话题。

多台车需要使用不同的节点名称、命名空间和TF坐标系，同时保持地图原点与时间一致。虚拟机能够通过SSH访问，不代表机器人一定能反向连接ROS；我们还需要检查局域网、\code{ROS_MASTER_URI}和\code{ROS_IP}。联调时应关闭不使用的速度发布程序，保证每台底盘只有一个主要控制源。
